  % ==================================================================
  \section{why tiny models?}
  % ==================================================================
	\begin{frame}{you (probably) have already used one}
    \begin{itemize}
      \item a chatbot, a translator, code completion.
      \item under each is a neural network: arrays of parameters and matrix
				multiplication.
      \item \textcolor{mute}{you could build one.}
    \end{itemize}
  \end{frame}

  \begin{frame}{frameworks are for production code}
    \hero{for learning, we will resort to minimalism.}
  \end{frame}

  \begin{frame}{the practice}
    \hero{build a \copperaccent{small, runnable} version yourself.}
    \subhero{mostly python and numpy and small enough to read fully.}
  \end{frame}

  \begin{frame}{what you trade}
    \begin{itemize}
      \item[\textcolor{cwTeal}{+}] visibility: every intermediate value is in front of you.
      \item[\textcolor{cwTeal}{+}] debuggability: a wrong value traces to one line.
      \item[\textcolor{cwTeal}{+}] intuition: the shapes show what each piece is
				for (puzzles, y’all!).
      \item[\textcolor{mute}{\textbar}] \textcolor{mute}{out of scope: scale, throughput, benchmark numbers.}
    \end{itemize}
  \end{frame}

  \begin{frame}{the setup}
    \begin{itemize}
      \item just \copperaccent{numpy arrays}, empasizing the shapes (they are
				most of the program).
      \item the notation in papers is shorthand for the array code.
    \end{itemize}
    \subhero{a trip to the terminal to build a transformer ourselves.}
  \end{frame}

  % ==================================================================
  \section{a tiny transformer, end to end}
  % one orienting slide, then the live build
  % ==================================================================
  \begin{frame}{what a transformer computes}
    \vfill
    \begin{columns}[T]
      \begin{column}{0.5\textwidth}
        {\ttfamily\large\color{cwTeal}%
          \renewcommand{\arraystretch}{1.7}%
          \begin{tabular}{@{}r@{~~}l@{}}
                  & tokens \\
            $\to$ & embeddings \\
            $\to$ & (attn + mlp) $\times N$ \\
            $\to$ & outputs \\
          \end{tabular}}
      \end{column}
      \begin{column}{0.46\textwidth}
        {\cwhankenbody\normalsize\color{cwInk}
          a stack of identical layers. each updates every position from all positions.\par}
        \vspace{1em}
        {\cwhankenbody\small\color{cwMuted}
          a token is a chunk of text, an embedding its vector. attention mixes
          across positions, the mlp transforms each position on its own.\par}
      \end{column}
    \end{columns}
    \vfill
  \end{frame}

  \begin{frame}{self-attention}
    \claim{a weighted average of all positions, with data-dependent weights.}
      {scores = Q K$^\top$ / $\sqrt{d_k}$ \\[0.4em] out = softmax(scores) V}
      {Q, K, V are linear projections of the input. scores is (T, T): how much each
       position draws from the others.}
  \end{frame}

  \begin{frame}[fragile]{softmax}
    \vfill
\begin{lstlisting}[language=Python]
def softmax(x):
    x = x - x.max(-1, keepdims=True)   # row max
    e = np.exp(x)
    return e / e.sum(-1, keepdims=True)
\end{lstlisting}
    \vskip6pt
    \subhero{softmax maps a vector to a probability distribution. subtracting the
    row max is the numerically stable form: identical result, no overflow in
    \texttt{exp}.}
  \end{frame}

  \begin{frame}{residual connections}
    \claim{the addition lets gradients reach every layer.}
      {x = x + attn(ln(x)) \\[0.4em] x = x + mlp(ln(x))}
      {the residual connection, from resnets (2015). deep stacks need it to train.}
  \end{frame}

  \begin{frame}{layer normalization}
    \claim{keeps activations from drifting as depth grows.}
      {y = $\gamma$ \,(x $-\ \mu$) / $\sqrt{\sigma^2 + \varepsilon}$\, $+\ \beta$}
      {$\mu, \sigma$: mean and variance over the feature axis, per position.
       $\gamma, \beta$: learned. $\varepsilon$: avoids dividing by zero.}
  \end{frame}

  \livecode{let's build one.}

  \begin{frame}[t]{a transformer block}
    \vskip0.09\paperheight
    \cwstat{$\sim$60}{lines of numpy, no framework}
    \vskip0.07\paperheight
    {\color{cwRule}\hrule height0.6pt}\vskip10pt
    {\cwhankenbody\color{cwFaint}\small attention, layer normalization, an mlp,
		residual connections, and the tests that pin them down. production models
		are this block, only wider and deeper. but: the operations are the same!}
  \end{frame}

  % ==================================================================
  \section{paper to prototype}
  % the transferable method
  % ==================================================================
  \begin{frame}{read an equation as an array program}
    \begin{itemize}
      \item every symbol is an array, every subscript is an axis.
      \item a $\sum_j$ is a matmul or a \texttt{.sum(axis=j)}.
      \item write the \copperaccent{shapes} in a comment before the code.
    \end{itemize}
    \subhero{if the shapes line up, you usually understand the mechanics of the formula.}
  \end{frame}

  \begin{frame}{catch bugs early with boring things}
    \begin{itemize}
      \item \copperaccent{shape asserts} at every boundary, not just the end.
      \item \copperaccent{unit tests for the math}: a known input, a hand-checked output.
      \item the \copperaccent{overfit test}: fit ten examples exactly. if you can't, it's broken.
    \end{itemize}
  \end{frame}

  % ==================================================================
  \section{multi-scale modelling (rgm)}
  % second architecture, shown from verified source
  % ==================================================================
  \begin{frame}{structure across scales}
    \hero{many signals carry structure at \copperaccent{several scales at once}:\\[0.3em]
    images, audio, text, physical fields.}
    \subhero{model the coarse structure first, then refine it.}
  \end{frame}

  \begin{frame}{coarse to fine}
    \claim{a resolution pyramid: model coarse structure first, detail after.}
      {down: (N) $\to$ (N/2) \\[0.3em] up: (N/2) $\to$ (N)}
      {coarse levels are small and global, fine levels carry details.}
  \end{frame}

  \begin{frame}{renormalization}
    \claim{apply the \emph{same} operator at every scale.}
      {x$_\ell$ = block(down(x$_{\ell+1}$)) \\[0.4em] for each level $\ell$}
      {borrowed from the renormalization group in physics: coarse-grain, model,
       repeat. \texttt{block} is the same weights at every level $\ell$.}
  \end{frame}

  \begin{frame}{generativity}
    \hero{run the pyramid the other way: start coarse,\\[0.3em]
    \copperaccent{invent detail} scale by scale.}
    \subhero{a learned model at each scale adds the next level of detail. image
		example: it produces new images, it doesn't just process the ones you give it.}
  \end{frame}

  % ==================================================================
  \section{state-space models}
  % third architecture, source not live
  % ==================================================================
  \begin{frame}{attention has a cost problem}
    \hero{it compares every pair of positions,\\[0.3em]
    so doubling the input \copperaccent{quadruples} the work.}
    \subhero{state-space models do sequence modelling in linear time: audio, dna, whole documents.}
  \end{frame}

  \begin{frame}{a state-space model}
    \hero{a loop with memory: walk the sequence,\\[0.3em]
    carry a running state, \copperaccent{update it each step}.}
    \subhero{\texttt{state = A\,state + B\,input}, then \texttt{output = C\,state}. like an rnn, but every step is linear.}
  \end{frame}

  \begin{frame}{one model with two forms}
    \hero{because the loop is linear, it is also\\[0.3em]
    a single \copperaccent{convolution} over the sequence.}
    \subhero{train as the convolution (whole sequence in parallel), run as the
		recurrence (one step, cheap memory). using the same weights!}
  \end{frame}

  \begin{frame}{attention also has a memory problem}
    \hero{to make the next token, it holds\\[0.3em]
    \copperaccent{every earlier token} in memory.}
    \subhero{that store grows with the sequence and never shrinks.}
  \end{frame}

  \begin{frame}{constant memory}
    \hero{a \copperaccent{fixed-size} state summarises the whole past,\\[0.3em]
    however long the sequence.}
    \subhero{and it's durable: save or reload it. state is a first-class object.}
  \end{frame}

  % ==================================================================
  \section{your own tiny model lab}
  % ==================================================================
  \begin{frame}{a reusable harness}
    \hero{we looked at three very different models.\\[0.3em]
    we can use the \copperaccent{same scaffold} for each.}
    \subhero{datasets, tests, math. build infra for your lab!}
  \end{frame}

  \begin{frame}{when to stop}
    \hero{the toy is for \copperaccent{understanding}.\\[0.3em]
    when you need scale, reach for pytorch, jax, w/e.}
  \end{frame}

  \begin{frame}{the takeaway}
    \hero{a paper is a specification.\\[0.3em]
    \copperaccent{implement the smallest version that runs}.}
		\subhero{try to stay true to the ideas only.}
  \end{frame}
